\documentclass{article}

\usepackage[utf8]{inputenc}     % pdfLaTeX
\usepackage[T1]{fontenc}
\usepackage[magyar]{babel}
\usepackage{amsmath}

% Set page size and margins
% Replace `letterpaper' with `a4paper' for UK/EU standard size
\usepackage[a4paper,top=2cm,bottom=2cm,left=3cm,right=3cm,marginparwidth=1.75cm]{geometry}
%for arrows
\usepackage{textcomp}

\title{A számnév}
\author{Lengyel Zoltán Péter\thanks{Csukás Ibolya tanítása alapján}}
\date{Legutóbb frissítve: 2026. Március 16.}

\begin{document}

\maketitle
\tableofcontents
\newpage

\section{A számnév fajtái}
\begin{itemize}
    \item Határozott számnevek:
    \[
\begin{array}{l l}
\left.
\begin{array}{l}
\text{Tőszámnév (pl. öt)} \\
\text{Törtszámnév (pl. ötöd)}
\end{array}
\right\}
& \text{Mennyiséget fejez ki}
\\[10pt]
\left.
\begin{array}{l}
\text{Sorszámnév (pl. ötödik)}\\
\end{array}
\right.\rightarrow
& \text{Sorrendiséget fejez ki}
\end{array}
\]
\item Határozatlan számnevek:
\begin{itemize}
    \item pl. sok, kevés, néhány
\end{itemize}
\end{itemize}


\section{A számnevek toldalékai}
\subsection{Képzők}
\begin{itemize}
    \item Sorszámnév képzője: -dik (ötödik)
    \item Törtszámnév képzője: -d (ötöd)
\end{itemize}
\subsection{Jelek}
\begin{itemize}
    \item A határozatlan számnevek fokozhatóak (kevés$\rightarrow$kevesebb$\rightarrow$legkevesebb)
\end{itemize}
\subsection{Ragok}
\begin{itemize}
    \item -an, -en
    \item -szor, -szer, -ször
\end{itemize}

\section{A számnevek helyesírása}
\begin{itemize}
    \item Kétezerig minden szám nevét egybe írjuk.
    \begin{itemize}
        \item pl. nyolcszázöt, ezerötszáznegyvennégy
    \end{itemize}
    \item Kétezer felett az összetett számneveket hármas számcsoportoknak megfelelően kötőjellel tagoljuk
    \begin{itemize}
        \item pl. kétezer-tizenegy, ötezer-háromszázhatvan
        \item de! kivételt képeznek az évszámok
    \end{itemize}
    \item Ha a sorszámneveket számjegyekkel írjuk, pontot teszünk utánuk
    \begin{itemize}
        \item A pontot megtartjuk, ha sorszámnevekhez toldalékot kapcsolunk
        \begin{itemize}
            \item pl. 6.-ban, 10.-be, 34.-en
        \end{itemize}
    \end{itemize}
    \item Év, hónap, nap után pontot teszünk
    \begin{itemize}
        \item Ha a dátum valamelyik elemét toldalékoljuk, az egyszerűsítés elve miatt elhagyjuk a pontot
        \begin{itemize}
            \item pl. 2020. január 1-jén, 1848 március 15-én
        \end{itemize}
    \end{itemize}
\end{itemize}

\end{document}